\documentclass[10pt,a4paper]{amsart}
\usepackage[margin=19mm]{geometry}
\usepackage{amsmath,amssymb,mathtools}
\usepackage[scr=rsfs,cal=euler]{mathalfa}
\usepackage{xfrac}
\usepackage[unicode,hidelinks,bookmarks=false]{hyperref}
\newcommand{\ie}{i.e.\ }
\newcommand{\cf}{cf.\ }
\newcommand{\resp}{respectively\ }
\newcommand{\Subsection}{\subsection}
\providecommand{\nobreakdash}{\nobreak\hbox{-}}
\setlength{\emergencystretch}{2em}
\multlinegap0pt
\usepackage{siunitx}
\usepackage{siunitx}
\usepackage{siunitx}
\usepackage{siunitx}
\usepackage{xcolor}
\usepackage{amssymb}
\usepackage{mathtools}
\usepackage{dsfont}
\usepackage{float}
\usepackage{enumerate}
\usepackage[shortlabels]{enumitem}
\usepackage{siunitx}
\usepackage{tikz}
\newtheorem{theorem}{Theorem}
\newcommand{\Gal}{\mathrm{Gal}}
\newcommand{\bQ}{\mathbb{Q}}
\title{Machines Learn Number Fields, But How? \\ The Case of Galois Groups}
\author{Kyu-Hwan Lee, Seewoo Lee}
\date{Source: arXiv:2508.06670v2 (CC BY 4.0)}
\begin{document}
\maketitle

\noindent\emph{Source and licence.} Machines Learn Number Fields, But How? \\ The Case of Galois Groups. Version arXiv:2508.06670v2 (CC BY 4.0), distributed by arXiv under the Creative Commons Attribution 4.0 International licence, http://creativecommons.org/licenses/by/4.0/, as recorded in the arXiv licence metadata for this exact version and shown on the abstract page https://arxiv.org/abs/2508.06670v2. Full licence notice: \texttt{licenses/arxiv-2508.06670-CC-BY-4.0.txt}.\par\smallskip

\noindent\textit{Editorial scope.} Counted unit: the article's theorem thm:cyclic\_degree\_lpower (source0.tex lines 986--992). The article's complete original proof of it is delivered with it (source0.tex lines 993--1007, one \texttt{proof} environment of 15 lines). No mathematics is written, repaired or re-derived: the statement and the proof are the article's own lines, in the article's own notation, and no retained line is substituted or re-flowed.  No further source-local context is needed: every cross-reference of every retained line resolves inside the two delivered spans.  Nothing was omitted from any retained span, so the package carries no omission record.  No retained line contains a citation, so the wrapper carries no bibliography and no bibliography entry.  No retained line contains a figure environment, an image-inclusion command or drawing code, so no image is delivered and the package declares no asset dependency.  Editorial formatting only, in addition: an AMS \texttt{amsart} preamble that restates the article's own declarations and macros used by the retained lines, namely the environments \texttt{theorem} on separate counters, together with the standard packages those lines need.  The retained environments print in this document as Theorem 1 in order of appearance, whereas the article prints the same items with the numbering of its own full document.  The complete native source of the article as distributed in the arXiv e-print for this exact version is bundled unchanged as \texttt{sources/182951/source0.tex}.
\begin{theorem}
    \label{thm:cyclic_degree_lpower}
    Let $\ell$ be a prime and $K / \bQ$ be a cyclic extension of degree $\ell^m$ for some $m \ge 2$, i.e. $G = \Gal(K / \bQ) \simeq C_{\ell^m}$.
     Then, for any $b \ge 1$ not divisible by $\ell^m$, the natural density of primes $p$ such that $a_{p^b}(K) = 0$ is positive. More precisely, if $\nu = v_\ell(b) < m$ is the $\ell$-adic valuation of $b$, then the density of primes with $a_{p^b}(K) = 0$ is $(\ell^{m - \nu} - 1) / \ell^{m - \nu}$.
    
    % Then the density of primes $p$ such that $a_{p^b}(K) = 0$ for any $b \in \mathbb Z_{\ge 1}$ not divisible by $\ell^m$ is positive.
\end{theorem}

\begin{proof}
    
    Consider unramified prime $p$ of decomposition type $(e, f, g) = (1, \ell^{t}, \ell^{m-t})$.
    The corresponding Euler factor is 
    \[
        \zeta_{K, p}(s) = (1 - p^{-\ell^{t}s})^{-\ell^{m - t}} = \sum_{k \ge 0} \binom{\ell^{m-t} + k - 1}{\ell^{m-t} - 1} p^{-\ell^{t}k s}
    \]
    and this shows $a_{p^b}(K) = 0$ if and only if $t > \nu$, i.e. $\ell^t \nmid b$.
    For each $t > \nu$, $\mathrm{Frob}_p$ has order $\ell^{t}$ and there are $\varphi(\ell^t) = \ell^{t - 1}(\ell - 1)$-many elements in $C_{\ell^m}$ of order $\ell^t$, so the density of prime $p$ with $a_{p^b}(K) = 0$ is
    \[
    \sum_{\nu < j \le m} \frac{\ell^{j-1}(\ell - 1)}{\ell^m} = \frac{\ell - 1}{\ell^m} \frac{\ell^{\nu}(\ell^{m - \nu} - 1)}{\ell - 1} = \frac{\ell^{m - \nu} - 1}{\ell^{m -\nu}}.
    \]
    Note that the natural density and Dirichlet density coincides in Chebotarev density theorem.
    
\end{proof}

\end{document}
